\scanpage{91}

\[
(28)\quad
\left\{\begin{aligned}
\widetilde{\Pi}_{0,3}=\Pi_{0,3}\cdot\{1,\tau\}&\longrightarrow\widetilde{\Pi}'_{0,3}\\
\tau_0,\tau_1,\tau_\infty&\longmapsto
\sigma'_0,\sigma'_1,\sigma'_\infty
\end{aligned}\right.
\]

This being so, we define an \uline{isomorphism} \deleted{\illegible{1}}
\[
(29)\quad
\begin{aligned}
\widetilde{\Pi}'_{0,3}&\xrightarrow{\sim}\mathfrak S_{0,3}\\
(\sigma'_0,\sigma'_1,\sigma'_\infty)
&\longrightarrow
\left(\deleted{\unsure{\sigma^{\prime}_1}}
\underset{\substack{\inserted{\text{over }\dot\tau'_\infty}\\\Vert\\\inserted{\dot\sigma_\infty\tau}}}{\tau'_\infty},
\underset{\substack{\inserted{\text{over }\dot\tau'_0}\\\Vert\\\inserted{\dot\sigma_0\tau}}}{\tau'_0},
\underset{\inserted{\text{over }\tau}}{\tau_\infty}\right)
\end{aligned}
\]
also denoted $r_0,r_1,r_\infty$ (fundamental \uline{reflections})

\margin{This isomorphism is associated with a canonical homomorphism
$\varphi:\widetilde{\Pi}_{0,3}=\Pi_{0,3}\cdot\{1,\tau\}\xrightarrow{\text{epimorphism}}\mathfrak S_{0,3}$
inducing a canonical homomorphism
$\varphi:\Pi_{0,3}\xrightarrow{\text{epimorphism}}\mathfrak S_{0,3}^{+}$
allowing us to identify $\mathfrak S_{0,3}$ and $\mathfrak S_{0,3}^{+}$
with the quotients of the left-hand~\unsure{sides|sides} by the relations
$\ell_1^2=\ell_\infty^3=1$ \ldots}

which induces an isomorphism
\[
(30)\quad
\left\{\begin{aligned}
\Pi'_{0,3}&\xrightarrow{\sim}\deleted{\unsure{\Pi}}\mathfrak S_{0,3}^{+}\\
\rho_s=\ell'_0&\longmapsto\varepsilon_0=\sigma_\infty\rho\\
\sigma=\ell'_1&\longmapsto\sigma_\infty\\
\rho_f=\ell'_\infty&\longmapsto\rho^{-1}
\end{aligned}\right.
\]

\margin{\uline{N.B.}
$\left\{\begin{aligned}
\tau_\infty\tau'_0&=\varepsilon_0=\ell'_0\\
\tau'_\infty\tau_\infty&=\sigma_\infty=\ell'_1\\
\tau'_0\tau'_\infty&=\rho^{-1}=\ell'_\infty
\end{aligned}\right.$
i.e.
$\begin{aligned}
r_\infty r_1&=\ell'_0\;(=\varepsilon_0)\\
r_0r_\infty&=\ell'_1\;(=\sigma_\infty)\\
r_1r_0&=\ell'_\infty\;(=\rho^{-1})
\end{aligned}$}

We thus find that $\mathfrak S_{0,3}$ has \deleted{\unsure{gen}} generating relations
\[
(31)\quad
{\tau'_\infty}^{2}={\tau'_0}^{2}=\tau_\infty^2
=\bigl(\underbrace{\tau'_\infty\tau_\infty}_{\sigma_\infty}\bigr)^2
=\bigl(\underbrace{\tau'_0\tau'_\infty}_{\rho^{-1}}\bigr)^3=1,
\]
and $\mathfrak S_{0,3}^{+}$ has generating relations
\[
(32)\quad \ell'_\infty \ell'_1\ell'_0={\ell'_1}^{2}={\ell'_\infty}^{3}=1
\]
or also simply, eliminating $\ell'_0$ and setting
$\ell'_1=\sigma_\infty$ and ${\ell'_\infty}^{-1}=\rho$
\[
(33)\quad\boxed{\sigma_\infty^2=\rho^3=1}
\]

For the record, we are going to express the notable elements of $\mathfrak S_{0,3}^{+}$ in terms of $\sigma_\infty,\rho$:

\scanpage{92}
\mspage{335}

\[
(34)\quad\left\{\begin{aligned}
\varepsilon_0&=\sigma_\infty\rho\ \deleted{\illegible{1}}&
\ell_0&=\sigma_\infty\rho\sigma_\infty\rho=(\sigma_\infty\rho)^2\\
\varepsilon_1&=\rho\sigma_\infty&
\ell_1&=\rho\sigma_\infty\rho\sigma_\infty=(\rho\sigma_\infty)^2\\
\varepsilon_\infty&=\rho^{-1}\sigma_\infty\rho^{-1}&
\ell_\infty&=(\rho^{-1}\sigma_\infty\rho^{-1})^2
=\rho^{-1}\sigma_\infty\rho^{-1}\sigma_\infty\rho^{-1}
\end{aligned}\right.,
\]
\[
(35)\quad\left\{\begin{aligned}
\sigma_0&=\rho^{\deleted{\unsure{2}}}\sigma_\infty\rho^{\deleted{\unsure{2}}-1}
\quad(=\deleted{\unsure{\rho}}\varepsilon_1\rho^{-1}=\rho\varepsilon_1^{-1})\\
\sigma_1&=\rho^{-1}\sigma_\infty\rho^{\deleted{\unsure{2}}}
\quad(=\rho^{-1}\varepsilon_0\deleted{\unsure{\rho}}=\varepsilon_\infty\rho^{-1})
\end{aligned}\right.
\]

and in $\mathfrak S_{0,3}^{-}$ we have, in terms of $\tau'_\infty,\tau'_0,\tau_\infty$,
the expressions for \deleted{$\unsure{\tau_0,\tau_1,\tau_\infty,\tau'_0,\tau'_1,\tau'_\infty}$}
\inserted{$\tau_0,\tau_1,(\tau_\infty,\tau'_0),\tau'_1,(\tau'_\infty)$}
\[
(36)\quad\left\{\begin{aligned}
\tau_0&=\deleted{\unsure{\tau_0\tau'_0\tau_\infty\tau'_\infty}}\rho\tau_\infty\rho^{-1}
=\tau'_\infty\tau'_0\tau_\infty\tau'_0\tau'_\infty
\quad(=r_0r_1r_\infty r_1r_0)\\
\tau_1&=\rho^{-1}\tau_\infty\rho
=\tau'_0\tau'_\infty\tau_\infty\tau'_\infty\tau'_0
\quad(=r_1r_0r_\infty r_0r_1)\\
\tau'_1&=\rho\tau'_0\rho^{-1}
=\tau'_\infty\deleted{\unsure{\tau'_0}}\tau'_0\tau'_\infty
\quad(=r_0r_1r_0)
\end{aligned}\right.
\]
and
\[
\left\{\begin{aligned}
\tau'_\infty&=r_0\\
\tau'_0&=r_1\\
\tau_\infty&=r_\infty.
\end{aligned}\right.
\]

\uline{IV}) Relation with
$\mathrm{Gl}(2,\mathbb Z)'=\mathrm{Gl}(2,\mathbb Z)/\mathord{\pm}1
\simeq\mathrm{GP}(1,\mathbb Z)$.

This group operates on the Poincaré half-plane
\[
D=\{z\in\mathbb C\mid\Im z>0\}
\]
by
\[
(37)\quad\left\{\begin{aligned}
u_g(z)&=\frac{az+b}{cz+d}&\text{if }ad-bc&=+1\\
u_g(z)&=\frac{a\bar z+b}{c\bar z+d}&\text{if }ad-bc&=-1
\end{aligned}\right.
\qquad
\begin{gathered}
g=\begin{pmatrix}a&b\\c&d\end{pmatrix}\\
a,b,c,d\in\mathbb Z\\
ad-bc\in\{\pm1\}
\end{gathered}
\]

We have in $D$ a hatched fundamental domain, which corresponds to the spherical
triangle $(-\bar j,0,\tfrac12)$ by the \unsure{homomorphism|homeomorphism}
of the natural \unsure{covering}, $i$ corresponding to $\tfrac12$,
$-\bar j$ to $-\bar j$, and the point at infinity to $0$.
\deleted{Thus \unsure{we find}}\inserted{The reflections}
with respect to the three sides of this triangle are given by
\[
(38)\quad
\deleted{\unsure{\tau_0}}\unsure{r_0|\tau_0}(z)=\frac1{\bar z},\qquad
\unsure{r_1|\tau_1}(z)=1-\bar z,\qquad
\unsure{r_\infty|\tau_\infty}(z)=-\bar z
\]
\figure{\begin{tikzpicture}[x=2.2cm,y=2.2cm,line cap=round,line join=round,font=\scriptsize]
\draw[line width=.4pt] (-1.55,0) .. controls (-.7,.04) and (.5,-.025) .. (1.85,.01);
\draw[line width=.4pt] (-1,0) arc[start angle=180,end angle=0,radius=1];
\draw[line width=.5pt] (0,-.48)--(0,2.85);
\draw[line width=.5pt] (.5,-.65)--(.5,2.78);
\begin{scope}
\clip (0,1) arc[start angle=90,end angle=60,radius=1]--(.5,2.78)--(0,2.85)--cycle;
\foreach \h in {.96,1.24,1.52,1.80,2.08,2.36,2.64}
  \draw[line width=.35pt] (0,\h)--(.5,{\h+.30});
\end{scope}
\draw[line width=.7pt] (0,1)--(.075,1)--(.075,1.075)--(0,1.075);
\fill (.5,.8660254) circle[radius=.025];
\node[inner sep=0pt] at (-1.04,-.14) {$-1$};
\node[inner sep=0pt] at (-.10,-.14) {$0$};
\node[inner sep=0pt] at (.43,-.18) {$\frac12$};
\node[inner sep=0pt] at (1.04,-.14) {$1$};
\node[inner sep=0pt] at (-.13,1.12) {$i$};
\node[inner sep=0pt] at (.64,.94) {$-\bar j$};
\node[inner sep=0pt,anchor=east,rotate=16] at (-.07,1.94) {$\dot r_\infty=\dot\tau_\infty$};
\node[inner sep=0pt,anchor=west] at (.56,2.25) {$\dot\tau'_0=\dot r_1$};
\node[inner sep=0pt,rotate=-63] at (.23,.63) {$\dot\tau'_\infty=\dot r_0$};
\end{tikzpicture}}

\deleted{\unsure{They determine}} and the pairwise products of these reflections

\scanpage{93}

are
\[
(39)\quad\left\{\begin{aligned}
\unsure{r_\infty r_1|\tau_\infty\tau_1}(z)&=\ell'_0(z)=z-1\;\deleted{\unsure{;}}\\\unsure{r_0r_\infty|\tau_0\tau_\infty}(z)&=\ell'_1(z)=-\frac1z\\
\unsure{r_1r_0|\tau_1\tau_0}(z)&=\ell'_\infty(z)=1-\frac1z=\frac{z-1}{z}
\end{aligned}\right.
\]

We thus find an isomorphism
\[
(40)\quad\left\{\begin{aligned}
\mathfrak S_{0,3}&\xrightarrow{\sim}\mathrm{Gl}(2,\mathbb Z)'\simeq
\mathrm{Gl}(2,\mathbb Z)/\mathord{\pm}1\simeq\mathrm{GP}(1,\mathbb Z)\\
r_0=\tau'_\infty&\longmapsto r_0=\begin{pmatrix}0&1\\1&0\end{pmatrix}\\
r_1=\tau'_0&\longmapsto r_1=\begin{pmatrix}-1&1\\0&1\end{pmatrix}\\
r_\infty=\tau_\infty&\longmapsto r_\infty=\begin{pmatrix}1&0\\0&-1\end{pmatrix}
\end{aligned}\right.
\]
\margin{it is better to change \unsure{some|of} signs \unsure{cf. Formulary|cf. formulary}}

and a presentation of $\mathrm{GP}(1,\mathbb Z)$ by
\[
(41)\quad r_0^2=r_1^2=r_\infty^2=(r_0r_\infty)^2=(r_1r_0)^3=1.
\]
This isomorphism induces an isomorphism
\[
(42)\quad\left\{\begin{aligned}
\mathfrak S_{0,3}^{+}&\xrightarrow{\sim}
\mathrm{Sl}(2,\mathbb Z)'=\mathrm{Sl}(2,\mathbb Z)/\mathord{\pm}1
\simeq\mathrm{GP}(1,\mathbb Z)^{+}\\
\ell'_1=\sigma_\infty&\longmapsto\begin{pmatrix}0&-1\\1&0\end{pmatrix}\\
\ell'_\infty=\rho^{-1}&\longmapsto\begin{pmatrix}1&-1\\1&0\end{pmatrix}
\end{aligned}\right.
\]
\[
(\text{hence}\quad\rho\longmapsto\begin{pmatrix}0&1\\-1&1\end{pmatrix}).
\]

We henceforth identify the elements of $\mathfrak S_{0,3}$ with their images in
$\unsure{\mathrm{Gl}|\mathrm{Sl}}(2,\mathbb Z)'\simeq\mathrm{GP}(1,\mathbb Z)$.
Let us note the following formulary \unsure{(where we use} that for a matrix
$u=\begin{pmatrix}a&b\\c&d\end{pmatrix}\in\mathrm{Sl}(2,\mathbb Z)$,
\deleted{\illegible{1}}
\deleted{\unsure{for} $\det u\in\mathbb Z^*=\{\pm1\}$ so $(\det u)^2=1$}, we have

\scanpage{94}
\mspage{336}

\[
(43)\quad u^{-1}=\begin{pmatrix}d&-b\\-c&a\end{pmatrix}\quad):
\]

\margin{N.B. \unsure{2.} \deleted{\unsure{If det}} $u\in\mathrm{Gl}(2,A)$,
$A$ an \unsure{arbitrary} ring, we have
$\det(u)\cdot u^{-1}=\begin{pmatrix}d&-b\\-c&a\end{pmatrix}$}

\deleted{\unsure{And}}.
\[
(44)\quad\left\{\begin{aligned}
\varepsilon_0=\sigma_\infty\rho&=\begin{pmatrix}1&-1\\0&1\end{pmatrix},&
\varepsilon_0^{-1}&=\begin{pmatrix}1&1\\0&1\end{pmatrix}\\
\varepsilon_1=\rho\sigma_\infty&=\begin{pmatrix}1&0\\1&1\end{pmatrix},&
\varepsilon_1^{-1}&=\begin{pmatrix}1&0\\-1&1\end{pmatrix}\\
\varepsilon_\infty=\rho^{-1}\sigma_\infty\rho^{-1}&=\begin{pmatrix}2&-1\\1&0\end{pmatrix},&
\varepsilon_\infty^{-1}&=\begin{pmatrix}0&1\\-1&2\end{pmatrix}\\
\ell_0=\varepsilon_0^2&=\begin{pmatrix}1&-2\\0&1\end{pmatrix}&
\sigma_0&=\begin{pmatrix}1&-1\\2&-1\end{pmatrix}\\
\ell_1=\varepsilon_1^2&=\begin{pmatrix}1&0\\2&1\end{pmatrix}&
\sigma_1&=\begin{pmatrix}1&-2\\1&-1\end{pmatrix}\\
\ell_\infty=\varepsilon_\infty^2&=\begin{pmatrix}3&-2\\2&-1\end{pmatrix}&
\sigma_\infty&=\begin{pmatrix}0&-1\\1&0\end{pmatrix}
\end{aligned}\right.
\]
\margin{Cf. ``formulary'' for a more \unsure{complete} list of relations and expressions.}
\[
(45)\quad\rho=\begin{pmatrix}0&1\\-1&1\end{pmatrix},\qquad
\unsure{\rho^{-1}|\rho^{-1\prime}}=\begin{pmatrix}1&-1\\1&0\end{pmatrix},
\]
\[
(46)\quad \ell_\infty \ell_1\ell_0=\begin{pmatrix}-1&0\\0&-1\end{pmatrix}
\]
the nontrivial element of the center of $\mathrm{Sl}(2,\mathbb Z)$ which
(fortunately!) gives $1$ in $\mathrm{Sl}(2,\mathbb Z)'$

The group $\Pi_{0,3}$ is recovered as the kernel of the homomorphism
\diagram{\begin{tikzcd}[ampersand replacement=\&]
\mathrm{Sl}(2,\mathbb Z)' \arrow[r] \& \mathrm{Sl}(2,\mathbb Z/2\mathbb Z) \arrow[d,phantom,"\simeq"{description,rotate=90}]\\
\& \mathfrak S_3
\end{tikzcd}}
the exact sequence
\[
(47)\quad
\diagram{\begin{tikzcd}[ampersand replacement=\&]
1 \arrow[r] \& \Pi_{0,3}\arrow[r,dash] \& \mathfrak S_{0,3}^{+}\arrow[r,dash] \& \mathcal{T}_{0,3}^{+\unsure{\prime}}\arrow[r,dash]\arrow[d,phantom,"\simeq"{description,rotate=90}] \& 1\\
\&\&\&\mathfrak S_3\&
\end{tikzcd}}
\]
is recovered as
\[
(48)\quad
\diagram{\begin{tikzcd}[ampersand replacement=\&]
1\arrow[r]\&\Pi(2)'\arrow[r]\arrow[d,phantom,"\simeq"{description,rotate=90}] \& \mathrm{Sl}(2,\mathbb Z)'\arrow[r] \& \mathrm{Sl}(2,\mathbb Z/2\mathbb Z)\arrow[r]\arrow[d,phantom,"\simeq"{description,rotate=90}] \& 1\\
\&\Pi_{0,3}^{\unsure{\beta}}\&\&\mathfrak S_3\&
\end{tikzcd}}
\]